\section{讨论：设计取舍、演进与教训}
\label{sec:discussion}

前面各章按模块讲清了“做了什么、怎么做”。本章换一个角度，讨论贯穿全系统的几个取舍：哪些工作该交给模型、哪些该交给程序；为什么用编码智能体做 harness；闸门与评分如何避免被“优化穿透”；哪些尝试失败了、为什么失败。我们也把系统三周多的演进史整理出来，因为每一次架构调整都是由具体失败驱动的，这些失败比最终设计本身更能说明设计为什么是现在的样子。

\subsection{模型与程序的分工边界}
\label{sec:disc-boundary}

“模型读、抄、写、判断；程序校验、计算、拦截”这条总原则在实践中需要一个可操作的判据：拿到一个新的子任务，怎么决定它归谁？我们最终使用的判据是两个问题：

\begin{enumerate}
  \item \textbf{结果能否被廉价地验证？}如果存在一个确定性程序，能以远低于完成任务本身的代价判断结果对错，那么任务可以交给模型，但必须配上这个验证器。抄录年报数字属于此类：抄录需要理解版式、区分合并与母公司报表，难以用规则穷尽；但验证“这个数字是否按数字边界出现在第 $p$ 页原文中”只需一次字符串扫描（第 \ref{sec:extract} 节）。
  \item \textbf{结果本身是否可由规则唯一确定？}如果答案由明确的公式唯一确定（同比增速、三表勾稽、折现），就不应交给模型，即便模型“多半能算对”。理由不只是准确率：可由规则确定的结果必须可复现、可做敏感性分析、可以被后续步骤当作常量引用，这三点采样生成都给不了。
\end{enumerate}

两个问题把任务分成四类（表 \ref{tab:disc-quadrant}）。值得注意的是右下格：既无法廉价验证、也无法由规则确定的任务（例如“这家公司的核心驱动是什么”“与市场的分歧在哪里”），系统没有假装能保证它们的正确性，而是交给多会话分工、复核编辑与盲评兜底，并在报告中把这类内容标为“观点”。这是系统诚实性的一部分：可验证的部分做到逐位可验证，不可验证的部分明确标出。

\begin{table}[h]
\centering
\caption{任务归属的两维判据。}
\label{tab:disc-quadrant}
\small
\begin{tabularx}{\linewidth}{@{}p{2.3cm}XX@{}}
\toprule
 & 结果可由规则唯一确定 & 结果不能由规则唯一确定 \\
\midrule
可廉价验证 & 程序直接计算（增速、占比、勾稽、DCF、评级）；模型不参与 & 模型执行 + 程序验证（抄录、经营数据表、来源摘句、研报断言拆分） \\
难以验证 & 程序计算并以 golden 快照、参照实现对拍（三表引擎、可比公司公式） & 模型执行 + 多会话交叉 + 盲评兜底，文中标为“观点”（投资主线、驱动推导、预期差） \\
\bottomrule
\end{tabularx}
\end{table}

\begin{tradeoff}[：让模型计算、程序复核]
一种折中方案是允许模型自己计算派生数，再由程序复算核对，不一致则拒收。它看起来兼顾了灵活性与正确性，我们没有采用，原因有三。第一，复核需要知道模型“想算的是什么”：模型写下“毛利率提升 2.3 个百分点”，程序必须先推断出这是哪两期、哪个口径的毛利率之差，推断本身就是一个易错的语义解析问题。第二，拒收后的重试循环会消耗大量轮数，\fab{模型在被拒收后倾向于删掉数字而不是改对，文稿因而变空}。第三，如果程序已经能复算，直接让程序算并把结果写进产物供模型引用，路径更短、失败模式更少。最终方案把“约数也不许心算”定为硬规则（第 \ref{sec:principles} 节），模型的任务从“算对”变成“找到并引用程序算好的数”，后者正是护栏可以机械检查的。
\end{tradeoff}

\subsection{系统演进：每一步由哪个失败驱动}
\label{sec:disc-evolution}

图 \ref{fig:disc-timeline} 按时间整理了系统的主要演进。从仓库提交历史看，系统大体经历了五个阶段。

\begin{figure}[t]
\centering
\begin{tikzpicture}[
  ev/.style={draw=nipsblue,fill=nipsblue!6,rounded corners=2pt,align=left,font=\scriptsize,text width=4.6cm,inner sep=3pt},
  dt/.style={font=\scriptsize\bfseries,text=nipsblue,anchor=east},
  dot/.style={circle,fill=nipsblue,inner sep=1.6pt}]
  \draw[thick,midgray] (0,0.3) -- (0,-12.1);
  \foreach \y/\d/\t/\side in {
    0/09-22/{初版：财报 / 估值 / 纠错端到端，依赖原 OpenAI 插件引擎}/r,
    -1.35/09-24/{自研替代原插件引擎，golden 快照约束口径一致}/l,
    -2.7/09-28/{Skill 版：Claude Code 直接读 PDF，程序定位、校验、计算；统一真值表并阻断未校验抽取}/r,
    -4.05/09-28/{EBIT 口径修正、三表驱动 DCF、季报单季拆分}/l,
    -5.4/10-02/{隔离工作区 + 分阶段闸门；多执行器（含国产模型）测评}/r,
    -6.75/10-03/{先搜集后分析；评分标准 v0、盲评、refine 闭环}/l,
    -8.1/10-04/{评分标准 v1（人工质检）；深度专题；真实年报跨页表格；评审校准}/r,
    -9.45/10-05/{评分标准 v2；反向 DCF；工作簿审计；出处分类；删除程序内 LLM 与其他执行器，只留一条路线}/l,
    -10.8/10-05/{渲染宽度与实测字号检查（修复整体缩放）}/r}{
    \node[dot] at (0,\y) {};
    \if\side r
      \node[ev,anchor=west] at (0.35,\y) {\t};
      \node[dt] at (-0.25,\y) {\d};
    \else
      \node[ev,anchor=east] at (-0.35,\y) {\t};
      \node[dt,anchor=west] at (0.25,\y) {\d};
    \fi
  }
\end{tikzpicture}
\caption{\sysname{} 的演进时间线（2026 年 9--10 月，依据仓库提交记录整理）。}
\label{fig:disc-timeline}
\end{figure}

\paragraph{阶段一：确定性引擎先行（9 月下旬）。}最早的版本是一条以规则抽取为主的确定性流水线，计算核心借用了一个开源插件的引擎。这一阶段确立了此后从未动摇的两件事：单一真值表，以及单位约定。随后我们用自研实现替换了外部引擎，并以 golden 快照逐项锁住计算口径（第 \ref{sec:valuation} 节）。这一步的代价是大量重复劳动，收益是对每一个计算口径都有完全的控制权，后面的 EBIT 口径修正、三表驱动 DCF 都依赖这一点。

\paragraph{阶段二：模型从“调用的 API”变成“工作的会话”（9 月 28 日前后）。}早期版本在程序内部调用模型 API 完成抽取，程序向模型发送页面文本、接收结构化 JSON。这种方式的问题在真实年报上暴露得很快：一份年报两三百页，程序必须事先决定把哪些页喂给模型，而“哪些页有用”恰恰需要理解版式；模型拿不到原始 PDF 的表格结构，也无法自己翻页核对。Skill 版把方向倒了过来：模型作为一个能读文件、能跑命令的会话直接读 PDF，程序退到“定位、校验、计算”的位置。这一转向是整个系统最重要的一次架构调整，“证据约束抽取”与“双通道不静默取舍”都是在此之后才真正可行的。

\paragraph{阶段三：从单会话到隔离工作区与闸门（10 月 2 日）。}单个会话写完整研报时，我们观察到三类问题：模型在仓库里能看到其他公司的分析框架并加以照抄；上下文被无关文件占满；模型会在后期阶段回头“修正”前期已经校验过的数字。隔离工作区与分阶段闸门同时解决了这三点：工作区里只有这一家公司的材料，前期产物一旦通过闸门，后续阶段改动它就会因新鲜度规则使闸门重新失效。

\paragraph{阶段四：先搜集后分析与评测闭环（10 月 3--4 日）。}第一批多会话研报在盲评中只得到 46～49 分（评分标准 v0），主要失分集中在行业与竞争章节：外部数字没有出处，或出处与数字对不上。把资料搜集提前为独立阶段、由研究助理会话完成并登记来源，是对这一失败的直接回应。同一时期加入了评分标准、盲评与 refine 闭环，评分标准随后依据人工质检修订为 v1，再把“数字前后矛盾”从模型判定改为程序判定形成 v2。

\paragraph{阶段五：收敛到一条路线（10 月 5 日）。}系统一度同时维护三种执行方式：程序内调用模型的工具循环、Claude Code 无头会话、以及另一家编码智能体执行器；还曾接入国产模型做对比测评。维护多条路线的成本在闸门与评测加入之后急剧上升：每加一道闸门、每改一次工作区说明，都要在三条路线上分别验证。我们最终删除了其他路线，只保留无头 Claude Code 一条，并在文档中说明了模型选择的理由（第 \ref{sec:disc-harness} 节）。

\begin{insight}[：演进史的共同模式]
回看这五个阶段，每一次调整的方向都相同：把一个原本依赖模型“自觉”的性质，改成由程序在边界上强制。照抄别家框架靠工作区隔离强制；不改已校验数字靠新鲜度规则强制；外部数字有出处靠来源登记与护栏强制；评分不泄露靠目录隔离强制。模型的能力没有被削弱，被削弱的是模型“违反约束而不被发现”的可能性。
\end{insight}

\subsection{为什么用编码智能体做 harness}
\label{sec:disc-harness}

本系统的每个模型角色都是一次无头 Claude Code 会话，而不是由我们自己的程序循环调用模型 API。这一选择有几方面的考虑。

\paragraph{能力面。}研报生产需要的能力——读 PDF、按页翻查、运行命令、读写文件、联网检索与抓取、在失败后自行重试——恰好是编码智能体已经打磨过的工具集。自建工具循环意味着重新实现文件系统沙箱、权限、上下文压缩、长任务中断恢复等基础设施，这些都不是本项目的贡献所在。

\paragraph{边界清晰。}以“会话 + 文件”为接口，模型与引擎之间只有一种交互方式（第 \ref{sec:arch} 节）。引擎因此可以保持纯粹的确定性，不含任何提示词或模型调用；所有与模型相关的配置集中在编排器生成的工作区说明与权限白名单里。这使得“引擎不调用模型”可以作为一条可测试的不变量存在。

\paragraph{可替换性。}会话的启动参数、模型名与工具白名单都由编排器传入，模型通过环境变量配置。理论上任何能以无头方式运行、支持文件读写与命令执行的编码智能体都可以替换当前的执行器。历史上我们确实接入过其他执行器并跑通过完整流程，删除它们是出于维护成本，不是架构不支持。

\begin{tradeoff}[：国产模型与 Claude]
大赛鼓励优先使用国产模型。我们曾用国产模型执行器完成过两家公司的研报流程并做对比，最终全部改用 Claude，理由是：（1）流水线对长程工具使用的稳定性要求很高，单个会话常常需要上百轮工具调用，执行器在这一点上的差异会直接表现为闸门通过率的差异；（2）维护多条执行路线使每次闸门或工作区规则的修改都要多次验证，在初赛时间约束下不可承受；（3）系统的正确性保证并不依赖于某个特定模型——数字正确由程序保证、来源由登记与护栏保证、评分由隔离的盲评保证——因此模型选择主要影响效率与文字质量，而非可信度。换回国产模型在架构上只需替换执行器与环境变量，我们将其列为决赛前的后续工作。
\end{tradeoff}

\begin{tradeoff}[：确定性编排器 vs 模型编排器]
多智能体框架常见的做法是由一个“经理”模型决定下一步调用谁。我们的编排器是一段确定性 Python：阶段顺序固定，是否进入下一阶段完全由闸门谓词决定，refine 的路由由问题单类型查表决定。我们放弃了模型编排的灵活性，换来三点：（1）流程可预测，同样的工作区状态总是触发同样的下一步，便于复现与调试；（2）闸门无法被“说服”，模型无法通过在总结里声称“已完成”来跳过检查；（3）编排器可以被单元测试覆盖，包括用模拟会话测试路由与回滚。代价是流程中的异常情况（例如外部网站长时间不可达）需要人工处理，编排器不会自己想办法绕过。
\end{tradeoff}

\subsection{闸门、评分与 Goodhart 定律}
\label{sec:disc-goodhart}

任何可度量的指标一旦成为优化目标，就会被优化穿透。流水线中有两类度量：闸门（程序检查，决定能否继续）与评分（盲评，决定质量）。我们对两类度量采取了相反的可见性策略：

\begin{itemize}
  \item \textbf{闸门完全可见。}闸门的条件写在工作区说明里，生成方清楚地知道需要满足什么。这是有意的：闸门检查的都是“必要条件”（数字可追溯、校验晚于修改、篇幅下限），被针对性满足正是我们想要的结果。被“优化穿透”的风险集中在数量型闸门上（字数、图表数、来源数），我们因此只把它们设为下限，并明确由盲评兜底质量。
  \item \textbf{评分完全不可见。}评分标准、分值与完整评测结果都不进入工作区，回传的问题单只含理由与证据。评分检查的是“质量”，一旦可见，生成方就会朝条目文本优化，例如在每节末尾机械地补一句“与一致预期的差异在于……”。
\end{itemize}

这一分法的底层逻辑是：对\textbf{必要条件}，可见性有利于满足；对\textbf{质量判断}，可见性导致表面满足。问题单之所以保留理由与证据，是因为“哪里不好、为什么”是改进所必需的信息，而“哪一条值几分”只对刷分有用。

\begin{failcase}[：数量下限的边界]
\fab{篇幅闸门引入后，早期版本出现过“深度研究”部分大段复述年报原文以凑足字数的情况。}我们没有通过提高闸门复杂度来应对（例如检测与年报的文本重合度），而是在主笔技能中写明深度的可检查要求（例如核心驱动至少两环定量推导、每个投资主线对应预测参数），并在盲评中设置对应条目。原因是任何新的数量检查都会被新的方式满足，而盲评条目对生成方不可见，更难被针对。
\end{failcase}

\subsection{误报的不对称代价贯穿全系统}
\label{sec:disc-fp}

“误报控制优先于召回”最初是为研报纠错核查制定的，但在系统的几乎每个检查点上都能看到同一取舍：

\begin{itemize}
  \item \textbf{抽取：}证据校验对找不到原文的数字直接拒收，宁可让真值表缺一项，也不让未经证实的数字进入下游。缺失可以被发现和补齐，错误数字一旦进入真值表会被估值、文稿、图表层层引用。
  \item \textbf{护栏：}容差按展示精度的半个末位设定，而不是一个宽松的相对容差，避免文稿中错误的数字“碰巧”匹配上某个产物数值而被放行。这里“误报”指的是护栏错放错误数字，其代价是错误进入交付物。
  \item \textbf{纠错：}错误类型只有在找到精确替代解释时才下结论，预测性表述判为“无法核实”。一个错报会让使用者怀疑全部核查结果。
  \item \textbf{评审：}盲评判出的问题要附原文证据，人工抽查关注的首要指标是精确率而非召回率。
\end{itemize}

这些检查点的共同点是：下游会信任上游的输出。在一条信任链上，每个环节的误报都会被放大，而漏报至少还有被下一环节捕获的机会。对信任链上游的检查偏向保守，是一种系统层面的风险分配，而不是各模块各自的偏好。

\subsection{不预设行业框架的代价与收益}
\label{sec:disc-noframework}

技能只给方法与约束，不给“白酒看吨价、电池看 GWh、航运看运价”这类行业模板。这一决定的收益是明确的：不同行业的公司得到的分析框架确实不同，模型会依据年报的分部披露与资料包推导驱动；模板不适用时也不会被硬套。代价同样明确：

\begin{enumerate}
  \item 框架质量依赖模型对行业的理解，对冷门行业的公司，推导出的驱动可能流于表面；
  \item 不同运行之间的框架差异更大，降低了同一公司多次运行结果的可比性；
  \item 盲评无法依据行业知识清单逐条检查，只能检查“是否给出了至少两环定量推导”这类与行业无关的结构性要求。
\end{enumerate}

我们接受这些代价，因为预设模板的失败模式更隐蔽：模板会让报告“看起来专业”，而读者很难察觉模板对这家公司并不适用。后续计划在拿到优秀研报样例后提炼“分析逻辑”而非“行业结论”，以改善第一点。

\subsection{没有成功的尝试}
\label{sec:disc-negative}

如实记录失败的尝试，有助于读者判断哪些设计是经过比较后的选择。

\paragraph{用正则拦截机器腔。}我们曾编写规则拦截“不是……而是……”“值得注意的是”“相互印证”等句式。\fab{规则上线后，模型很快改用同义表达，文风并没有改善，反而在被拦截后出现语句不通的修补。}最终删除了规则，改为由不参与写作的复核编辑会话改写，数字锁保证改写不碰数字（第 \ref{sec:guards} 节）。

\paragraph{程序内工具循环。}早期版本让程序维护一个“提问—调用工具—回答”的循环，由模型决定调用哪个分析函数。\fab{它在样例年报上工作良好，在真实年报上频繁因上下文长度、工具参数格式和中断恢复出问题。}与其继续打磨这一循环，不如直接使用成熟的编码智能体，这一路线最终被删除。

\paragraph{整节删除类缺陷的评审检出。}在评审校准中，我们用程序造出“删掉一个必要章节”的劣化稿，期望评审能判出。实际检出率偏低。分析发现主要原因在劣化算子本身：删除一节正文后，摘要、图表与附录中仍保留了相关内容，评审据此判定“已覆盖”。这说明劣化算子需要跨全文同步删除，也提醒我们，评审校准的结论受限于劣化的彻底程度（第 \ref{sec:eval} 节）。

\paragraph{单次大会话写完整研报。}在多会话流水线之前，我们尝试过让一个会话完成从抽取到排版的全部工作。\fab{会话在后期常因上下文压缩丢失前期约束，例如忘记某个数字已被拒收，或在排版阶段修改正文。}拆分为五个角色后，每个会话的上下文只包含本阶段所需的材料，这类问题基本消失。

\subsection{可迁移性}
\label{sec:disc-transfer}

尽管本系统针对 A 股财报，核心机制并不依赖金融领域：
\begin{itemize}
  \item \textbf{证据约束抽取}适用于任何“从长文档中抄出结构化事实”的任务，如合同条款、临床试验报告、审计底稿；
  \item \textbf{展示精度容差}适用于任何需要判断“文本中的数字是否与真值一致”的场景；
  \item \textbf{数字锁}适用于任何“允许模型改写、但不允许改动事实”的编辑任务；
  \item \textbf{闸门可见、评分不可见}的分离适用于任何用模型评审驱动模型生成的闭环。
\end{itemize}
领域相关的部分集中在指标 schema、勾稽规则、估值引擎与技能说明中。迁移到港股或美股需要替换报表版式处理与会计口径映射；迁移到其他文档密集型领域则需要替换真值表的计算层。
